% Part: turing-machines
% Chapter: undecidability
% Section: introduction

\documentclass[../../../include/open-logic-section]{subfiles}

\begin{document}

\olfileid{tur}{und}{int}
\olsection{پېژندنه}

ښايي دا څرګنده وبرېښي چې هره تابع، ان هره حسابي تابع، محاسبه کېدونکې
نۀ شي کېدے۔ تابعې دومره ډېرې دي او چلند ئې دومره پېچلے دے۔ مثلاً هغه
تابعې چې د راډيو‌اکټيفو ذرو له تجزيې، يا له بل ګډوډ يا تصادفي چلند
څخه تعريفېږي۔ فرض کړئ له يوه ټاکلي وخت څخه يوه-ثانيه‌يي واټنونه شمېرو،
او تابع~$f(n)$ د کاينات د هغو ذرو د شمېر په توګه تعريفوو چې له هغې
لومړنۍ شېبې وروسته په~$n$-م يوه-ثانيه‌يي واټن کښې تجزيه کېږي۔ دا د
داسې تابعې غوندې ښکاري چې هېڅکله ئې د محاسبې هيله نۀ شو کولے۔

خو دا چې د داسې تابعو د محاسبې لار تصور نۀ شو کولے يوه خبره ده، او
دا چې په رښتيا ثابته کړو چې دوي نۀ محاسبه کېږي، بله خبره ده۔ په اصل
کښې ان هغه تابعې چې بېخي ناهيلې کوونکې پېچلې ښکاري، په مجردې معنا
محاسبه کېدونکې کېدے شي۔ مثلاً فرض کړئ کاينات د وخت له پلوه متناهي دے
---لکه ځينې کيهان‌پوهنيزې تيورۍ چې وړاندوينه کوي، يوه ورځ به په ډېر
لرې راتلونکي کښې کاينات په يوه ټکي کښې راټول شي۔ بيا له هغې لومړنۍ
شېبې څخه يوازې د متناهي (خو بېخي لوے) شمېر ثانيو دپاره~$f(n)$ تعريف
شوې وي۔ او هره تابع چې يوازې د متناهي شمېر ننوتونو دپاره تعريف شوې
وي محاسبه کېدونکې ده: وتونه په يوه لوے جدول کښې لست کولے، يا ئې د
ټيورينګ ماشين د حالت بدلون په يوه ډېر لوے ډياګرام کښې کوډولے شو۔

موږ ډېر ځله د هغو تابعو ځانګړو حالتونو ته پام کوو چې ارزښتونه ئې د
هو/نه پوښتنو ځوابونه ورکوي۔ مثلاً پوښتنه ``ايا~$n$ لومړنی عدد دے؟''
له دې تابع سره تړلې ده:
\[
\fn{isprime}(n) = \begin{cases}
  1 & \text{که $n$ لومړنی وي}\\
  0 & \text{که نۀ وي۔}
  \end{cases}
\]
وايو يوه هو/نه پوښتنه په \emph{اغېزمنه توګه پرېکړه کېدے} شي، که ور
سره تړلې د~$1/0$ ارزښت لرونکې تابع په اغېزمنه توګه محاسبه کېدونکې وي۔

د دې د رياضياتي ثبوت دپاره چې داسې تابعې شته چې په اغېزمنه توګه نۀ
شي محاسبه کېدے، يا داسې مسئلې شته چې په اغېزمنه توګه نۀ شي پرېکړه
کېدے، اړينه ده چې د محاسبې يو ټاکلی مدل وټاکو، او وښيو چې داسې تابعې
شته چې دا ئې نۀ شي محاسبه کولے يا داسې مسئلې شته چې دا ئې نۀ شي
پرېکړه کولے۔ مثلاً ښودلے شو چې هره تابع د ټيورينګ ماشينونو په وسيله
نۀ شي محاسبه کېدے او هره مسئله د ټيورينګ ماشينونو په وسيله نۀ شي
پرېکړه کېدے۔ بيا د چرچ--ټيورينګ اصل کارولے شو، څو پايله واخلو چې نۀ
يوازې ټيورينګ ماشينونه د هرې تابع د محاسبې دپاره پوره ځواکمن نۀ دي،
بلکې هېڅ اغېزمنه کړنلاره دا کار نۀ شي کولے۔
\textbf{د ژباړې د نسخې سمون (OLCMP-048):} په سرچينه کښې د ``په
اغېزمنه توګه پرېکړه کېدے'' په غونډله کښې مرستندوی فعل پاتې شوے و؛
دلته بشپړ ګرامري جوړښت راوړل شوے دے او معنا نۀ ده بدله شوې۔

د داسې منفي پايلو د ثبوت کلي دا حقيقت دے چې موږ پخپله ټيورينګ
ماشينونو ته عددونه ټاکلے شو۔ تر ټولو اسانه لار ئې دا ده چې دوي وشمېرو:
ښايي د ټيورينګ ماشينونو او د هغو د پروګرامونو د ليکلو يوه ټاکلې لار
وټاکو، او بيا ئې په منظم ډول لست کړو۔ کله چې ووينو دا کار کېدے شي،
نو د ټيورينګ-نامحاسبه کېدونکو تابعو شتون د شمېر کچې له ساده دلايلو
څخه راځي: له~$\Nat$ څخه~$\Nat$ ته د تابعو سټ (په اصل کښې ان يوازې
له~$\Nat$ څخه~$\{0, 1\}$ ته) !!{nonenumerable} دے، خو څرنګه چې ټول
ټيورينګ ماشينونه شمېرلے شو، د ټيورينګ-محاسبه کېدونکو تابعو سټ يوازې
!!{denumerable} دے۔

موږ داسې \emph{ټاکلې} تابعې او مسئلې هم تعريفولے شو چې په ترتيب سره
ئې نۀ محاسبه کېدل او ناپرېکړتيا ثابتولے شو۔ يوه داسې مسئله د
\emph{درېدو مسئلې} په نوم ده۔ د ټيورينګ ماشين لارښوونې په لست کولو
هغه په متناهي ډول بيانېدے شي۔ د ټيورينګ ماشين دا ډول بيان، يعنې د
ټيورينګ ماشين پروګرام، البته بل ټيورينګ ماشين ته د ننوت په توګه
کارېدے شي۔ نو هغه ټيورينګ ماشينونه څېړلے شو چې د نورو ټيورينګ
ماشينونو په اړه پوښتنې پرېکړه کوي۔ يوه په زړه پورې پوښتنه دا ده:
``ايا ورکړل شوے ټيورينګ ماشين به پر ننوت~$n$ له پيل وروسته بالاخره
ودرېږي؟'' ښه به وے که داسې ټيورينګ ماشين موجود وے چې دا پوښتنه ئې
پرېکړه کولے: دا د کيفيت-کنټرول يو ټيورينګ ماشين وګڼئ چې ډاډ ورکوي
ټيورينګ ماشينونه په نامتناهي کړيو او داسې نورو کښې نۀ بندېږي۔ په زړه
پورې حقيقت، چې ټيورينګ ثابت کړ، دا دے چې داسې ټيورينګ ماشين نشي
موجودېدے۔ هېڅ يو ټيورينګ ماشين نشته چې د يوه ټيورينګ ماشين~$M$ له
بيان او يوه عدد~$n$ څخه جوړ ننوت باندې پيل شي او تل له~$1$ يا~$0$
وت سره د دې له مخې ودرېږي چې~$M$ به پر ننوت~$n$ درېدلے و که نه۔
\textbf{د ژباړې د نسخې سمون (OLCMP-049):} سرچينې په وروستۍ پوښتنه
کښې د اېم له نوم وروسته د ``ماشين'' کلمه په تکرار راوړې وه؛ دلته
زائده تکرار لرې شوے او منطقي شرط هماغسې ساتل شوے دے۔

کله چې د ټاکلو ناپرېکړه کېدونکو مسئلو بېلګې ولرو، د نورو مسئلو د
ناپرېکړتيا د ښودلو دپاره ئې هم کارولے شو۔ مثلاً يوه مشهوره
ناپرېکړه کېدونکې مسئله دا پوښتنه ده: ``ايا د لومړۍ درجې
!!{formula}~$!A$ معتبر دے؟'' داسې ټيورينګ ماشين نشته چې د لومړۍ درجې
!!{formula}~$!A$ د ننوت په توګه ورکړل شي او هرومرو له~$1$ يا~$0$ وت
سره د دې له مخې ودرېږي چې~$!A$ معتبر دے که نه۔ په تاريخي ډول د دې
مسئلې د اغېزمن حل د کړنلارې د موندلو پوښتنې ته يوازې ``د'' پرېکړې
مسئله ويل کېده؛ نو وايو چې د پرېکړې مسئله ناحل کېدونکې ده۔ ټيورينګ
او چرچ دا پايله نږدې په يوه وخت کښې په خپلواکه توګه ثابته کړه، ځکه
نو دې ته د چرچ--ټيورينګ قضيه هم ويل کېږي۔

\end{document}
